% ============================================================
\section{Related Work}
\label{sec:related}
% ============================================================

Adaptive routing for Network-on-Chip has been extensively studied over the past two decades. This section reviews relevant literature across three categories: classic heuristic adaptive routing, deep reinforcement learning approaches, and graph neural network methods for network optimization.

\subsection{Classic Adaptive Routing}

Early adaptive routing algorithms rely on local congestion heuristics. DyAD~\cite{dyad2004} dynamically switches between deterministic and adaptive routing modes based on congestion thresholds. When local buffer occupancy exceeds a threshold, the router switches to adaptive mode to explore alternative paths. Regional Congestion Awareness (RCA)~\cite{ebrahimi2019} extends this concept by propagating congestion information across multiple hops, providing a wider view of network state. However, these methods use fixed, hand-crafted thresholds that cannot adapt to changing traffic patterns. Moreover, they rely on local or limited-region information, missing the global topology structure.

Deadlock-free adaptive routing theory was established by Duato~\cite{duato1996,duato1993}, who proved necessary and sufficient conditions for deadlock-free routing in wormhole networks. The Odd-Even turn model~\cite{odd2000} and West-First routing~\cite{glass1992} provide deadlock-free adaptive routing by restricting certain turns. While these methods are theoretically sound, they trade routing flexibility for deadlock freedom and cannot exploit global congestion information.

\subsection{Deep Reinforcement Learning for NoC Routing}

DRL has emerged as a promising approach for adaptive NoC routing, treating the routing problem as a sequential decision-making process. MAAR~\cite{wang2022maar} proposes multi-agent DRL where each router hosts an independent Q-network, trained with centralized experience replay. Their approach achieves 20-35\% latency improvement over XY routing across multiple traffic patterns. However, MAAR uses an MLP encoder that flattens router state into a fixed vector, discarding the spatial relationships between neighboring routers.

DeepNR~\cite{deepnr2022} applies deep Q-networks to NoC routing with a focus on energy efficiency. Their approach achieves 18\% energy reduction while maintaining performance. GARNN~\cite{garnn2023} incorporates graph attention mechanisms but stops short of providing a full NoC-specific GNN architecture for routing decisions. Q-learning based approaches~\cite{choudhary2026} have been proposed for fault-tolerant routing, using RL to bypass faulty links.

Chan et al.~\cite{chan2023lare} proposed LARE (Linear Approximate RL), which replaces traditional Q-tables with a linear function approximator to achieve adaptive routing with significantly reduced hardware overhead compared to full Q-learning. Their approach maintains comparable performance to Q-learning while reducing storage requirements by an order of magnitude. Khan et al.~\cite{khan2023qrasp} introduced Q-RASP, a regional congestion-aware RL routing policy that shares path experience across packets sharing similar routes, achieving 18.3\% average packet latency improvement on 4$\times$4 and 8$\times$8 meshes. Wang et al.~\cite{wang2024drlar} developed DRLAR, a deep RL adaptive routing framework that formulates routing as an agent considering multiple network state features and system-level feedback, demonstrating effectiveness against state-of-the-art adaptive routing algorithms.

The key limitation across DRL-based approaches is their use of fully-connected (MLP) encoders. An MLP processes router state as a flat vector of features, without any mechanism to capture the graph structure of the NoC topology. As we demonstrate in Section~\ref{sec:problem}, topology structure significantly affects congestion patterns, suggesting that topology-aware representations could improve routing decisions.

\subsection{Graph Neural Networks for Network Routing}

GNNs have been applied to routing in wide-area and software-defined networks (SDN). The GROM framework~\cite{grom2024} combines DRL with GNNs to generate routing policies that generalize to unseen network topologies without retraining, addressing a key limitation of topology-specific models. G-Routing~\cite{g-routing2023} combines GNNs with DRL for routing in communication networks, demonstrating that GNN-based policies generalize across different network topologies. DRL-GNN~\cite{drl-gnn2019} uses GNN encoders for routing optimization, showing that message-passing architectures capture graph structure better than MLPs. However, these approaches target networks with millisecond-scale latency requirements, not the nanosecond-scale constraints of on-chip routing.

For NoC-specific applications, GraphNoC~\cite{graphnoc2024} applies GNNs for application-specific FPGA NoC performance prediction, while PathGNN~\cite{pathgnn2025} proposes a path-based GNN architecture for robust and resilient distributed routing that quickly adapts to unexpected network conditions. DQN-FTR~\cite{dqnftr2025} uses Deep Q-Networks for fault-tolerant routing in mesh NoCs, demonstrating RL's potential for handling dynamic faults alongside congestion.

More recently, zero-shot generalization has emerged as a key research direction for ML-based routing. The ability to transfer knowledge from a small, easily-trained network to a larger deployment-scale network without retraining can substantially reduce training costs. Our work contributes to this direction by demonstrating zero-shot weight transfer from a 4$\times$4 GNN model to an 8$\times$8 mesh.

For chip design, GNNs have been used for NoC topology optimization. NOCTOPUS~\cite{noctopus2026} predicts NoC performance metrics using GNNs, achieving high accuracy for latency and throughput prediction. Circuit-GNN~\cite{mirhoseini2021} uses graph representations for chip placement optimization, demonstrating that GNNs can learn useful representations of circuit structure. GraphNoC~\cite{graphnoc2024} applies GNNs for application-specific FPGA NoC performance prediction, showing that GNN embeddings correlate with post-implementation metrics. However, these works use GNNs for prediction rather than routing decisions.

Inference latency is the critical barrier to GNN-based on-chip routing. A single GNN forward pass on a 64-node graph requires approximately $10^6$ cycles on a CPU~\cite{xiao2024}, while NoC routing decisions must complete within $5-10$ cycles~\cite{wang2022maar}. This approximately five-order-of-magnitude gap necessitates architectural innovations such as periodic policy optimization~\cite{periodic2025} or dedicated GNN accelerators~\cite{xiao2024}.

\subsection{Quantitative Comparison with State-of-the-Art}

Table~\ref{tab:sota-comparison} provides a detailed comparative analysis of GNNocRoute against representative ML-based NoC routing works, covering encoder architecture, inference overhead, and key capabilities.

\begin{table}[h]
\centering
\caption{Comprehensive comparison with state-of-the-art ML-based NoC routing. ``N/C'' denotes not considered or not applicable.}
\label{tab:sota-comparison}
\resizebox{\columnwidth}{!}{%
\begin{tabular}{lcccccc}
\toprule
\textbf{Criterion} & \textbf{GNNocRoute} & \textbf{MAAR} & \textbf{DeepNR} & \textbf{GARNN} & \textbf{DQN-FTR} & \textbf{LARE} \\
 & \textbf{(Ours)} & \cite{wang2022maar} & \cite{deepnr2022} & \cite{garnn2023} & \cite{dqnftr2025} & \cite{chan2023lare} \\
\midrule
\textbf{Encoder} & GATv2 & MLP & MLP & GAT & MLP & Linear \\
\textbf{Topology-aware} & Yes ($r$=0.978) & No & No & Partial & No & No \\
\textbf{Training} & Offline & Online DQN & Online DQN & Offline & Online DQN & Online SARSA \\
\textbf{Per-hop inference} & 1 cycle & $N$ cycles & $N$ cycles & $10^3$ cycles & $N$ cycles & 2--3 cycles \\
\textbf{Area overhead} & 4--64 KB & 50--200 KB & 50--200 KB & MB$+$ & 50--200 KB & Few KB \\
\textbf{Latency gain vs. XY} & 25--45\% & 20--35\% & $\sim$15--20\% & N/C & 18--30\% & Comparable to QL \\
\textbf{Energy reduction} & 22--30\% & N/R & 18\% & N/C & N/R & N/R \\
\textbf{Fault tolerance} & Yes (retrain) & No & No & No & Yes & No \\
\textbf{Zero-shot (4$\times$4$\to$8$\times$8)} & 99.78\% & No & No & N/A & No & No \\
\textbf{Deadlock freedom} & Turn-model & Restricted turns & Escape VC & Restricted turns & Fault-aware & Restricted turns \\
\bottomrule
\end{tabular}%
}
%}\resizebox
\end{table}

\paragraph{Per-hop inference cost.}
The critical advantage of GNNocRoute lies in its inference efficiency. All DRL-based approaches that make online per-hop routing decisions (MAAR, DeepNR, DQN-FTR) require a full forward pass through a neural network at each router hop, imposing $5$--$20$~cycles of additional latency. LARE~\cite{chan2023lare} reduces this to $2$--$3$ cycles via linear approximation but sacrifices representation capacity. GNN-based methods (GARNN, G-Routing, GROM) are even costlier, requiring $10^3$--$10^4$ cycles for a full graph pass---impractical for NoC's $5$--$10$~cycle budget.

GNNocRoute decouples policy computation from inference: the GATv2-based DRL agent recomputes routing tables every $T$ cycles (periodic optimization), while each per-hop decision is a single-cycle table lookup. This achieves the inference latency of heuristic routing (1~cycle) with the routing quality of learned policies, making GNNocRoute the first practical GNN-based routing method that satisfies NoC timing constraints.

\paragraph{Topology awareness.}
MLP-based methods (MAAR, DeepNR, DQN-FTR) flatten router state into a fixed-length vector, discarding spatial relationships between neighboring routers. This is a fundamental architectural limitation: the encoder cannot distinguish congestion patterns that differ only in their topological arrangement. GARNN~\cite{garnn2023} uses graph attention but targets wide-area networks and lacks NoC-specific features. Our GATv2 encoder achieves $r=0.978$ correlation with betweenness centrality, demonstrating genuine topology awareness. This structure-aware representation directly contributes to the 25--45\% latency improvement over XY routing---outperforming MAAR's 20--35\% despite identical DRL training framework, confirming that the GNN encoder provides superior state representations.

\paragraph{Zero-shot generalization and fault tolerance.}
A distinguishing capability of GNNocRoute is its $99.78\%$ zero-shot generalization from a 4$\times$4 trained model to an 8$\times$8 mesh. No existing DRL-based NoC routing method achieves this---all require per-topology retraining (MAAR, DeepNR, DQN-FTR). GNN-based WAN routing methods (G-Routing, GROM) claim generalization but operate under $10^3\times$ more relaxed timing constraints.

Similarly, GNNocRoute supports retraining for fault tolerance: when a router fails, the GNN is retrained on the modified topology and the new routing table is distributed. This is a natural consequence of the offline/periodic paradigm, whereas online methods require complex reward engineering to handle faults during exploration.

\paragraph{Summary.}
GNNocRoute achieves comparable or better routing performance than state-of-the-art DRL approaches (MAAR, DeepNR) while requiring only $0.3$--$2\%$ of their per-hop inference cost. The GATv2 encoder provides genuine topology awareness absent in MLP-based methods, and the periodic optimization framework enables zero-shot generalization and fault tolerance as free byproducts. This combination of accuracy, efficiency, and flexibility positions GNNocRoute as the first practical GNN-based routing solution for NoC systems.

\subsection{Position of This Work}

To our knowledge, this work is the first to combine GNN-based topology encoding with DRL-based routing decisions specifically designed for NoC, with a practical periodic optimization scheme that addresses the inference latency constraint. Table~\ref{tab:related} positions our work relative to existing approaches.

\begin{table}[h]
\centering
\caption{Comparison with representative works}
\label{tab:related}
\resizebox{\columnwidth}{!}{%
\begin{tabular}{lccccl}
\toprule
\textbf{Work} & \textbf{Encoder} & \textbf{Decision} & \textbf{Objective} & \textbf{Latency-aware} \\
\midrule
DyAD~\cite{dyad2004} & Threshold & Local heuristic & Latency & Yes \\
RCA~\cite{ebrahimi2019} & Regional & Heuristic & Latency & Yes \\
MAAR~\cite{wang2022maar} & MLP & DRL & Latency & Yes \\
DeepNR~\cite{deepnr2022} & MLP & DQN & Latency+Energy & Yes \\
G-Routing~\cite{g-routing2023} & GNN & DRL & Latency & No (WAN) \\
GNNocRoute (Ours) & \textbf{GNN} & \textbf{DRL} & \textbf{Latency+Energy} & \textbf{Yes} \\
\bottomrule
\end{tabular}%
}
\end{table}
